\subsection{The drafts}
\label{sec:9.4.1}
Two bills in the current Congress answer the first question. The Abolish ICE Act would cut off federal funds for the functions assigned to ICE's director under the 2002 Act, rescind the unobligated balances of what had been made available to its director, and abolish the agency ninety days after enactment \autocite{hr7123}. The Melt ICE Act works through the agency's powers rather than its name \autocite{hr7190}:

\begin{itemize}
\item Everyone held in DHS detention is released on recognizance within six months.
\item The detention authorities of the Immigration and Nationality Act are repealed, including all of section 236.
\item Section 287 is rewritten without the 287(g) delegation to state and local police.
\item The two provisions that forbid state and local governments to restrict their officials' communication with immigration authorities, 8 U.S.C. 1373 and 1644, are repealed.
\item Detention and monitoring contracts end within two years, after which no federal funds may be used for detention or monitoring programs at all.
\item Ankle monitors come off within six months.
\item Within two years, funding ends for information-sharing partnerships between DHS and state or local law enforcement agencies that identify or target noncitizens.
\item ICE's operations account may not pay for civil arrests, detention, removal or charging documents.
\item An HHS grant program funds services for people affected by enforcement. Only organizations that have never taken part in immigration or law enforcement may deliver them, they may not surveil the people who use them, and they may not pass those people's personal information to any federal entity.
\end{itemize}

Two further drafts reach one entry of the schedule. The ICE Out of Our Faces Act, S. 3779, would bar ICE and CBP from acquiring or using facial recognition and other biometric identification systems, require deletion of the data collected for them, and let individuals and state attorneys general sue for violations \autocite{Jayapal2026FacesAct}. The Facial Recognition and Biometric Technology Moratorium Act, reintroduced in September 2026, would bar every federal entity from using facial recognition and other biometric technologies until Congress lifts the bar. It would also make federal grants to state and local agencies conditional on their adopting moratoria of their own \autocite{Jayapal2026Moratorium}.

This Congress passed the Secure America Act 214–212 and 52–47 on party lines \autocite{GT2026SecureAmerica}, and none of these bills has a path to enactment in it. They exist as drafts, written before anyone needs them, so that enacting one is a matter of votes and not of drafting.
